\subsection{LIE ALGEBRAS OF CERTAIN LINEAR GROUPS}

Let \(\mathbf{E}\) be a complete normable space. Then \(\mathcal{L}(\mathbf{E})\) is a unital complete normable algebra and \(\mathbf{GL}(\mathbf{E})\) is a Lie group. By the Corollary to Proposition 33, no. 9, if \(\mathrm{T}_{1}(\mathbf{GL}(\mathbf{E}))\) is canonically identified with \(\mathcal{L}(\mathbf{E})\) , the Lie algebra structure on \(\mathrm{L}(\mathbf{GL}(\mathbf{E}))\) is given by the bracket \((x,y)\mapsto xy - yx\) of two elements of \(\mathcal{L}(\mathbf{E})\) . In particular, \(\mathrm{L}(\mathbf{GL}(n,\mathbf{K}))\) is canonically identified with \(\mathrm{gl}(n,\mathbf{K})\) (Chapter I, §1, no. 2).

PROPOSITION 35. Let E be a finite-dimensional vector space. Let \(\phi\) be the morphism \(g \mapsto \det g\) of the Lie group \(\mathbf{GL}(E)\) into the Lie group \(\mathbf{K}^*\) . The mapping \(\mathrm{L}(\phi)\) of \(\mathcal{L}(E)\) into \(\mathbf{K}\) is the mapping \(x \mapsto \mathrm{Tr}x\) . The kernel \(\mathbf{SL}(E)\) of \(\phi\) is a Lie subgroup of \(\mathbf{GL}(E)\) with Lie algebra \(\mathfrak{sl}(E)\) .

We choose a norm and a basis of E. The expansion of the determinant proves that

\[
\det (1 + u) \in 1 + \operatorname{Tr} u + o (\| u \|)
\]

when \(u\) tends to 0 in \(\mathcal{L}(\mathrm{E})\) . Hence, using Proposition 34, no. 9, for \(x \in \mathcal{L}(\mathrm{E}) = \mathrm{L}(\mathbf{GL}(\mathrm{E}))\) :

\[
\mathrm{L} (\phi) (x) = \langle x, \phi \rangle = \operatorname{Tr} x.
\]

It follows that \(\phi\) is a submersion. Therefore, \(\operatorname{Ker} \phi = \mathbf{SL}(E)\) is a Lie subgroup of \(\mathbf{GL}(E)\) whose Lie algebra is \(\operatorname{Ker} L(\phi) = \mathfrak{sl}(E)\) .

Let \(\mathbf{E}_1, \ldots, \mathbf{E}_n\) be complete normable spaces and \(\mathbf{E}\) their direct sum. Every \(x \in \mathcal{L}(\mathbf{E})\) can be represented by a matrix \((x_{ij})_{1 \leqslant i, j \leqslant n}\) , where \(x_{ij} \in \mathcal{L}(\mathbf{E}_i, \mathbf{E}_j)\) .

PROPOSITION 36. Let \(\mathbf{I}\) be a subset of \(\{1,2,\ldots,n\}\) and \(\mathbf{G}\) the subgroup of \(\mathbf{GL}(\mathbf{E})\) consisting of the \(g = (g_{ij})_{1\leqslant i,j\leqslant n}\in \mathbf{GL}(\mathbf{E})\) such that \(g_{ij} = 0\) for \(i < j\) and \(g_{ii} = 1\) for \(i\in \mathbf{I}\) . Then \(\mathbf{G}\) is a Lie subgroup of \(\mathbf{GL}(\mathbf{E})\) and \(\mathbf{L}(\mathbf{G})\) is the set of \(x = (x_{ij})_{1\leqslant i,j\leqslant n}\in \mathcal{L}(\mathbf{E})\) such that \(x_{ij} = 0\) for \(i < j\) and \(x_{ii} = 0\) for \(i\in \mathbf{I}\) . Let \(\mathbf{S}\) be the set of \((x_{ij})\in \mathcal{L}(\mathbf{E})\) such that \(x_{ij} = 0\) for \(i < j\) and \(x_{ii} = 0\) for \(i\in \mathbf{I}\) . Then \(\mathbf{G}\) is the intersection of \(\mathbf{GL}(\mathbf{E})\) and the affine subspace \(1 + \mathbf{S}\) of \(\mathcal{L}(\mathbf{E})\) . Hence \(\mathbf{G}\) is a submanifold of \(\mathbf{GL}(\mathbf{E})\) and the tangent space to \(\mathbf{G}\) at 1 is identified with \(\mathbf{S}\) .

In particular, in \(\mathbf{GL}(n,\mathbf{K})\) , the total lower triangular subgroup and the lower strict triangular subgroup, defined as in Integration, Chapter VII, § 3, no. 3, are Lie subgroups with Lie algebras \(t(n,\mathbf{K})\) and \(n(n,\mathbf{K})\) (Chapter I, § 1, no. 2).

PROPOSITION 37. Let A be a complete normable unital associative algebra and \(x \mapsto x^t\) a continuous linear mapping of A into A such that \((x^t)^t = x\) and \((xy)^t = y^t x^t\) for all \(x, y\) in A. Suppose that K is of characteristic \(\neq 2\) . Let G be the subgroup of A* consisting of the \(x \in A\) such that \(xx^t = x^t x = 1\) . Then G is a Lie subgroup of A* and L(G) is the set of \(y \in A\) such that \(y^t = -y\) .

Let \(S\) (resp. \(S'\) ) be the set of \(y \in A\) such that \(y = y^t\) (resp. \(y = -y^t\) ). Then \(S, S'\) are closed vector subspaces of \(A\) . The formula

\[
y = \frac {1}{2} (y + y ^ {t}) + \frac {1}{2} (y - y ^ {t})
\]

proves that A is the topological direct sum of S and \(S'\) . Let f be the mapping of A into S defined by \(f(x) = xx^{t}\) . This mapping is analytic. For all \(y \in A\) , \(f(1 + y) = 1 + y + y^{t} + yy^{t}\) ; choose a norm on A compatible with its algebra structure; then

\[
f (1 + y) \in 1 + y + y ^ {t} + o (\| y \|) \quad \text { for   } y \text {   tending   to   } 0.
\]

Thus, \(\mathrm{T}_{1}(f)(y)=y+y^{t}\) , so that f is a submersion at 1. Therefore, there exists an open neighbourhood U of 1 in A such that \(U\cap G\) is a submanifold of U. Hence (§ 1, no. 3, Proposition 6) G is a Lie subgroup of \(A^{*}\) . Moreover, \(\mathrm{L}(G)=\mathrm{T}_{e}(G)=\mathrm{Ker}\;\mathrm{T}_{1}(f)\) .

COROLLARY 1. Suppose that K is of characteristic \(\neq 2\) . Let E be a finite-dimensional vector space over K and \(\phi\) a non-degenerate symmetric (resp. alternating) bilinear form on E. For all \(u \in \mathcal{L}(\mathrm{E})\) , let \(u^*\) be the adjoint of \(u\) relative to \(\phi\) . Let G be the orthogonal (resp. symplectic) group of \(\phi\) . Then G is a Lie subgroup of \(\mathbf{GL}(\mathrm{E})\) and \(\mathrm{L}(\mathrm{G})\) is the set of \(x \in \mathcal{L}(\mathrm{E})\) such that \(x^* = -x\) .

We apply Proposition 37 with A = L(E) and \(x^{t} = x^{*}\) .

Remark. Let B be a basis of E and J the matrix of \(\phi\) with respect to B. Then \(\mathrm{L}(G)\) is the set of elements of \(\mathcal{L}(\mathrm{E})\) whose matrix X with respect to B satisfies the equation

\[
{ } ^ { t } X = - J X J ^ { - 1 } .
\]

This follows from Algebra, Chapter IX, § 1, formula (50).

COROLLARY 2. Let E be a complex (resp. real) Hilbert space and U the unitary group of E. Then U is a real subgroup of \(\mathbf{GL}(\mathrm{E})\) and \(\mathbf{L}(\mathbf{U})\) is the set of \(x \in \mathcal{L}(\mathrm{E})\) such that \(x^{*} = -x\) .

We apply Proposition 37 with \(\mathbf{A} = \mathcal{L}(\mathbf{E})\) considered as an algebra over \(\mathbf{R}\) and \(x^t = x^*\) .

COROLLARY 3. Let E be a finite-dimensional complex vector space, \(\phi\) a non-degenerate Hermitian sesquilinear form on E and U the unitary group of \(\phi\) . Then U is a real Lie subgroup of \(\mathbf{GL}(\mathbf{E})\) and \(\mathbf{L}(\mathbf{U})\) is the set of \(x \in \mathcal{L}(\mathbf{E})\) such that ix is Hermitian.

When \(E \neq \{0\}\) , U is not a Lie subgroup of the complex Lie group \(\mathbf{GL}(E)\) , for \(\mathbf{L}(U)\) is not a complex vector subspace of \(\mathcal{L}(E)\) .

